% Section 10
\section{Appendix B: Effective Field Theory Structure}
\label{sec:eft}

The bilocal construction has a natural effective 
field theory interpretation. The right wedge 
$\mathcal{R}$ is the \emph{system}, with energy 
scale $E_{\text{sys}}$; the full Minkowski space 
$\mathcal{R} \cup \mathcal{L}$ is the 
\emph{environment}, with energy scale 
$E_{\text{env}} = E_{\text{sys}} + E_\mathcal{L}$. 
The Rindler horizon is the EFT cutoff surface: 
it is the boundary at which the system's 
description and the full Minkowski description 
become incompatible, and the left wedge 
$\mathcal{L}$ is precisely the degrees of freedom 
that the system has integrated out.

\subsection{The GHY Term as Effective Action}
\label{sec:eft_ghy}

We now show that the GHY boundary term arises 
as the effective action for $\mathcal{H}_R$ after 
integrating out $\mathcal{H}_L$ in the path 
integral, elevating the identification 
$S_{GHY}|_\mathcal{H} = S_{\text{matter}}$ from 
Section~\ref{sec:ppwave} from an algebraic 
observation to a structural result.

\paragraph{Setup.}
The full Euclidean path integral for gravity and 
matter decomposes over the two wedges as
%
\begin{equation}
Z = \int \mathcal{D}g_R\,\mathcal{D}\phi_R
    \int \mathcal{D}g_L\,\mathcal{D}\phi_L\;
    e^{-S_R - S_L - S_{\mathcal{H}}},
\end{equation}
%
where $S_{\mathcal{H}}$ is the boundary term at 
the horizon coupling the two wedges. Integrating 
out the left wedge defines the effective action:
%
\begin{equation}
e^{-S_{\text{eff}}[g_R,\phi_R]}
= \int \mathcal{D}g_L\,\mathcal{D}\phi_L\;
  e^{-S_L - S_{\mathcal{H}}}.
\label{eq:eff_action_def}
\end{equation}

\paragraph{Origin of the horizon coupling.}
Integration by parts applied to any kinetic term 
produces a boundary contribution at $\mathcal{H}$. 
For the gravitational field this is the 
Gibbons--Hawking--York term~\cite{gibbons1977action}
%
\begin{equation}
S_{\mathcal{H}} = S_{GHY}\big|_{\mathcal{H}} 
= \frac{1}{8\pi G}\int_{\mathcal{H}} 
  d^2x_\perp\,d\lambda\;\theta,
\end{equation}
%
which couples the two wedge geometries through 
their expansion $\theta$ at the horizon, precisely 
as a scalar kinetic term couples $\phi_L$ and 
$\phi_R$ through their normal derivatives.

\paragraph{Saddle point evaluation.}
The classical solution in the left wedge is flat 
Minkowski space perturbed by the AS pp-wave of the 
left-moving photon. Aichelburg--Sexl pp-waves are 
Ricci flat ($R_{\mu\nu} = 0$) in vacuum, so the 
bulk Einstein--Hilbert action vanishes on-shell:
%
\begin{equation}
S_L\big|_{\text{on-shell}} 
= \frac{1}{16\pi G}\int_L d^4x\sqrt{g}\,R = 0.
\end{equation}

In a standard EFT path integral, integrating out 
the left wedge would produce two contributions: 
the gravitational boundary term $S_{GHY}$ (from 
integration by parts of the bulk kinetic term) 
and the matter action $S_{\text{matter}}^L$ 
(from the left-moving photon). Treated as 
independent, these would give 
$e^{-S_{GHY} - S_{\text{matter}}^L}$, 
and the effective action would be 
$S_R + S_{GHY} + S_{\text{matter}}^L = 
S_R + 2\hbar\omega$.

The consistency check of 
Section~\ref{sec:ppwave} resolves this: 
$S_{GHY}$ and $S_{\text{matter}}^L$ are not 
independent contributions but two descriptions 
of the same physical event, one photon 
crossing the horizon, counted geometrically 
and informationally. Adding them separately 
double-counts that event. The correct horizon 
coupling is therefore $S_{\mathcal{H}} = 
S_{GHY}$ alone, with $S_{\text{matter}}^L$ 
already subsumed into the geometric boundary 
term. At the saddle point:
%
\begin{equation}
\left.
\int \mathcal{D}g_L\,\mathcal{D}\phi_L\;
e^{-S_L - S_{\mathcal{H}}}
\right|_{\text{saddle}}
= e^{-S_{GHY}|_{\mathcal{H}}},
\end{equation}
%
where $S_L|_{\text{on-shell}} = 0$ by 
Ricci-flatness, and the effective action for 
the right wedge is

Since $S_{GHY}|_{\mathcal{H}} = 
S_{\text{matter}}^L = \sum_i\hbar\omega_i$ 
by equation~\eqref{eq:GHY_micro}, the effective 
action for the right wedge is
%
\begin{equation}
\boxed{
S_{\text{eff}}[g_R,\phi_R]
= S_R[g_R,\phi_R] + S_{GHY}\big|_{\mathcal{H}}.
}
\label{eq:Seff}
\end{equation}
%
The GHY term is the contribution to the right 
wedge effective action from integrating out the 
left wedge at the classical level. The bulk left 
wedge action vanishes by Ricci-flatness of the 
AS solution; the entire left wedge effect is 
encoded in the boundary term at the horizon.

\paragraph{Quantum corrections.}
Fluctuations of $\phi_L$ around the saddle 
generate corrections to $S_{\text{eff}}$ 
suppressed by $\ell_P^2/A$ relative to the 
classical term. These are the standard 
entanglement entropy contributions studied by 
Susskind--Uglum~\cite{susskind1994black} and 
Callan--Wilczek~\cite{callan1994geometric}, 
whose UV divergences renormalize Newton's 
constant $G$. Equation~\eqref{eq:Seff} is the 
leading classical contribution.

\subsection{Partition Function and S-Matrix}
\label{sec:eft_partition}

The thermal partition function of the right wedge
%
\begin{equation}
Z = \mathrm{Tr}_{\mathcal{H}_R}\,
    e^{-H_R/k_BT_H}
\end{equation}
%
is obtained from the full Minkowski vacuum by 
tracing out $\mathcal{H}_L$: integrating out the 
left wedge in the path integral yields the thermal 
density matrix for $\mathcal{H}_R$, and $\ln Z$ 
generates the connected correlation functions of 
the right wedge EFT. The S-matrix of the full 
bilocal construction connects the past state at 
$\mathcal{P}$ to the future sink $\mathcal{F}$:
%
\begin{equation}
S = \langle\mathcal{F}|\mathcal{P}\rangle
  = \mathrm{Tr}\,e^{-iHt/\hbar},
\end{equation}
%
which under Wick rotation $\tau = it$ becomes 
$Z = \mathrm{Tr}\,e^{-H\tau/\hbar}$. The Carnot 
cycle of Section~\ref{sec:carnot} is therefore 
the thermodynamic avatar of this S-matrix element: 
the real-time bilocal construction and its 
thermodynamic description are related by Wick 
rotation, and each microstate~\eqref{eq:microstate_vec} 
contributes one term to $Z$ rather than the full 
thermal trace.

\subsection{Horizon Complementarity}
\label{sec:eft_complementarity}

The identification of $\mathcal{H}_R$ as system 
and $\mathcal{H}_L \otimes \mathcal{H}_R$ as 
environment is the operational content of horizon 
complementarity~\cite{susskind1993stretched}: the 
right wedge observer sees a thermal state at $T_H$, 
while the full Minkowski observer sees a pure 
entangled state. These are the EFT and UV 
descriptions of the same microstate. The spatial 
bit, the answer to ``who observed $H_+$?'', 
is visible to $\mathcal{R}$ as a classical record 
but its quantum origin in the entangled pair is 
accessible only to the full Minkowski observer who 
holds both wedges. Complementarity is not an 
additional postulate here but an operational 
consequence of the EFT structure: the horizon is 
the surface at which the single-wedge and 
two-wedge descriptions become incompatible, 
consistently with the equivalence list of the 
following subsection.

\subsection{EFT Validity and the Planck Scale}
\label{sec:eft_validity}

The construction defines two natural length scales. 
The system scale is the Compton wavelength of the 
detector
%
\begin{equation}
\lambda_C = \frac{\hbar}{m_\mathcal{R} c},
\end{equation}
%
the scale below which the single-particle 
description of $\mathcal{R}$ breaks down. The 
environment scale is the proper horizon distance
%
\begin{equation}
d_\mathcal{H} = \frac{c^2}{a},
\end{equation}
%
the scale at which the single-bit description 
saturates. The EFT is valid when these scales are 
separated, $\lambda_C \ll d_\mathcal{H}$; it 
breaks down when they coincide. The Planck length
%
\begin{equation}
\ell_P = \sqrt{\frac{\hbar G}{c^3}} 
       = \sqrt{\lambda_C \cdot r_s}
\label{eq:planck_geometric_mean}
\end{equation}
%
is the geometric mean of the Compton wavelength 
and the Schwarzschild radius $r_s = 2Gm/c^2$: 
it is the unique length scale built from both 
$\hbar$ and $G$ and is the scale at which 
$\lambda_C = r_s$, i.e.\ where the system scale 
and the gravitational scale coincide. The Planck 
scale is therefore the fixed point of the EFT: 
the unique scale at which the quantum and 
gravitational descriptions become equally valid 
and the EFT has no regime of applicability.

The Carnot efficiency of Section~\ref{sec:carnot}
measures the separation between these scales. 
With $T_H = \hbar a/2\pi k_Bc$ and $T_C = 
\hbar\kappa/2\pi k_Bc$, the efficiency
%
\begin{equation}
\eta = 1 - \frac{T_C}{T_H} 
     = 1 - \frac{\kappa}{a}
     = \frac{m_\mathcal{L}}
            {m_\mathcal{R} + m_\mathcal{L}}
\label{eq:eta_mass}
\end{equation}
%
is the fraction of the total mass residing in 
the environment, the degrees of freedom 
integrated out. The following conditions are 
simultaneously equivalent:

\begin{prop}[EFT Validity]
\label{prop:eft}
The following are equivalent for the bilocal 
construction:
\begin{enumerate}
    \item $\kappa < a$ \quad 
    (mirror surface gravity less than global 
    Rindler acceleration)
    \item $E_{\mathrm{sys}} < E_{\mathrm{env}}$ 
    \quad (system energy less than environment)
    \item $\eta > 0$ \quad 
    (Carnot cycle extracts positive work)
    \item $\lambda_C < d_\mathcal{H}$ \quad 
    (Compton scale inside horizon distance)
    \item $\ell_P / \lambda_C \ll 1$ \quad
    (equivalently $m_\mathcal{R} \gg m_P$, 
    the detector mass is super-Planckian)
    \item $\mathcal{R}$ is a proper subsystem of 
    its environment.
\end{enumerate}
The boundary case where all inequalities become 
equalities is simultaneously the horizon 
condition, the Planck condition 
$\lambda_C = r_s = \ell_P$, and the 
thermodynamic equilibrium condition $\eta = 0$. 
The EFT breaks down exactly at the horizon, 
which is not a failure of the framework but 
its natural boundary condition.
\end{prop}
